\chapter{Multivariate Gaussian Distributions, Maximum Likelihood, GMMs, EM Algorithm \& Concentration Inequalities}
\label{chap:week12}

% ----------------------------------------------------------------------
% PART 1: EXECUTIVE ROADMAP & LEARNING OBJECTIVES
% ----------------------------------------------------------------------
\section{Executive Roadmap \& Learning Objectives}

The multivariate normal (Gaussian) distribution is the cornerstone of statistical machine learning. By virtue of the Central Limit Theorem, aggregated noise fluctuations across physical systems converge toward Gaussianity. However, real-world data often exhibits multimodality and latent cluster structure that a single unimodal Gaussian cannot capture. \textbf{Gaussian Mixture Models (GMMs)} resolve this by modeling data as a convex combination of Gaussians, estimated iteratively via the **Expectation-Maximization (EM) algorithm**. Finally, **Concentration Inequalities** provide non-asymptotic probabilistic bounds governing sample sizes and generalization error.

\subsection{Prerequisites \& Scaffolding}
Mastery of this chapter synthesizes positive definite matrices, Jensen's inequality, and probability theory as organized in Table~\ref{tab:w12_prereqs}.

\begin{table}[htbp]
\centering
\small
\begin{tabularx}{\textwidth}{lXX}
\toprule
\textbf{Foundational Area} & \textbf{Core Concepts Required} & \textbf{Role in Machine Learning Foundations} \\
\midrule
Linear Algebra & Matrix inverse $\Sigma^{-1}$, determinant $|\Sigma|$, Spectral Theorem $\Sigma = Q \Lambda Q^T$ & Defining Mahalanobis distance ellipsoids and normalizing constants \\
Convex Analysis & Jensen's Inequality $f(\E[X]) \le \E[f(X)]$, concave logarithms & Deriving the Evidence Lower Bound (ELBO) for the EM algorithm \\
Probability & Conditional density, Bayes' rule, marginalization, variance & Formulating posterior responsibilities and mixture likelihoods \\
\bottomrule
\end{tabularx}
\caption{Prerequisite Scaffolding for Chapter 12.}
\label{tab:w12_prereqs}
\end{table}

\subsection{Core Competencies}
By the conclusion of this chapter, students will be able to:
\begin{itemize}
    \item Evaluate multivariate Gaussian densities and compute Mahalanobis distance metrics.
    \item Derive the Maximum Likelihood Estimators (MLE) for multivariate Gaussian mean $\boldsymbol{\mu}$ and covariance $\Sigma$.
    \item Formulate Gaussian Mixture Models with latent categorical cluster assignments.
    \item Derive the complete Expectation-Maximization (EM) algorithm from first principles via the Evidence Lower Bound (ELBO).
    \item State and apply Markov's, Chebyshev's, and Hoeffding's concentration inequalities to establish rigorous sample complexity bounds.
\end{itemize}

\subsection{Notation \& Mathematical Conventions}
\begin{itemize}
    \item $\mathcal{N}(\mathbf{x}; \boldsymbol{\mu}, \Sigma)$: The $d$-dimensional multivariate Gaussian probability density.
    \item $D_M(\mathbf{x}, \boldsymbol{\mu}) \triangleq \sqrt{(\mathbf{x} - \boldsymbol{\mu})^T \Sigma^{-1} (\mathbf{x} - \boldsymbol{\mu})}$: The Mahalanobis distance.
    \item $K$: The number of latent mixture components.
    \item $\pi_k \triangleq P(z = k)$: Mixture weight for component $k$ ($\sum_{k=1}^K \pi_k = 1, \pi_k \ge 0$).
    \item $\gamma_{ik} \triangleq P(z_i = k \mid \mathbf{x}_i)$: The posterior responsibility of component $k$ for sample $i$.
    \item $\text{ELBO}(q, \boldsymbol{\theta})$: The Evidence Lower Bound on the marginal log-likelihood.
    \item $\bar{X}_N \triangleq \frac{1}{N}\sum_{i=1}^N X_i$: The empirical sample mean.
\end{itemize}

% ----------------------------------------------------------------------
% PART 2: MATHEMATICAL & COMPUTATIONAL FOUNDATIONS
% ----------------------------------------------------------------------
\section{Mathematical \& Computational Foundations}

\subsection{The Multivariate Gaussian Distribution}

\begin{definitionbox}{Multivariate Normal Density}
A continuous random vector $\mathbf{X} \in \mathbb{R}^d$ follows a \textbf{Multivariate Gaussian distribution} $\mathbf{X} \sim \mathcal{N}(\boldsymbol{\mu}, \Sigma)$ with mean $\boldsymbol{\mu} \in \mathbb{R}^d$ and symmetric positive definite covariance matrix $\Sigma \succ 0$ if its PDF is:
\begin{equation}
\mathcal{N}(\mathbf{x}; \boldsymbol{\mu}, \Sigma) \triangleq \frac{1}{(2\pi)^{d/2} |\Sigma|^{1/2}} \exp\left( -\frac{1}{2}(\mathbf{x} - \boldsymbol{\mu})^T \Sigma^{-1} (\mathbf{x} - \boldsymbol{\mu}) \right)
\label{eq:mvn_def}
\end{equation}
\end{definitionbox}

\begin{theorem}[Geometric Contour Ellipsoids]
The level sets (contours of equal probability density) of a multivariate Gaussian are concentric ellipsoids characterized by constant Mahalanobis distance:
\begin{equation}
(\mathbf{x} - \boldsymbol{\mu})^T \Sigma^{-1} (\mathbf{x} - \boldsymbol{\mu}) = c^2
\label{eq:mahalanobis_ellipsoid}
\end{equation}
By the Spectral Theorem, $\Sigma = Q \Lambda Q^T$. The principal axes of the ellipsoid align with the orthogonal eigenvectors $\mathbf{q}_i$, with semi-axis lengths proportional to $\sqrt{\lambda_i}$.
\end{theorem}

\subsection{Maximum Likelihood Estimation for Gaussian Parameters}
Given i.i.d. observations $S = \{\mathbf{x}_1, \dots, \mathbf{x}_N\} \subset \mathbb{R}^d$, the log-likelihood function is:
\begin{equation}
\ell(\boldsymbol{\mu}, \Sigma) = -\frac{N d}{2}\log(2\pi) - \frac{N}{2}\log|\Sigma| - \frac{1}{2}\sum_{i=1}^N (\mathbf{x}_i - \boldsymbol{\mu})^T \Sigma^{-1} (\mathbf{x}_i - \boldsymbol{\mu})
\label{eq:mle_loglik}
\end{equation}

\begin{theorem}[Gaussian Maximum Likelihood Estimators]
The Maximum Likelihood Estimators for $\boldsymbol{\mu}$ and $\Sigma$ are:
\begin{align}
\hat{\boldsymbol{\mu}}_{\text{MLE}} &= \frac{1}{N} \sum_{i=1}^N \mathbf{x}_i \label{eq:mle_mu} \\
\hat{\Sigma}_{\text{MLE}} &= \frac{1}{N} \sum_{i=1}^N (\mathbf{x}_i - \hat{\boldsymbol{\mu}}_{\text{MLE}})(\mathbf{x}_i - \hat{\boldsymbol{\mu}}_{\text{MLE}})^T \label{eq:mle_sigma}
\end{align}
While $\hat{\boldsymbol{\mu}}_{\text{MLE}}$ is unbiased ($\E[\hat{\boldsymbol{\mu}}] = \boldsymbol{\mu}$), $\hat{\Sigma}_{\text{MLE}}$ is biased: $\E[\hat{\Sigma}_{\text{MLE}}] = \frac{N-1}{N}\Sigma$. The unbiased sample covariance incorporates Bessel's correction factor $\frac{N}{N-1}$.
\end{theorem}

\subsection{Gaussian Mixture Models (GMM) and the EM Algorithm}
A Gaussian Mixture Model models complex multimodal densities as a convex combination of $K$ distinct Gaussian components:
\begin{equation}
p(\mathbf{x}; \boldsymbol{\theta}) = \sum_{k=1}^K \pi_k \mathcal{N}(\mathbf{x}; \boldsymbol{\mu}_k, \Sigma_k)
\label{eq:gmm_density}
\end{equation}
where $\boldsymbol{\theta} = \{\pi_k, \boldsymbol{\mu}_k, \Sigma_k\}_{k=1}^K$.
Direct maximization of the marginal log-likelihood $\ell(\boldsymbol{\theta}) = \sum_{i=1}^N \log\left( \sum_{k=1}^K \pi_k \mathcal{N}(\mathbf{x}_i; \boldsymbol{\mu}_k, \Sigma_k) \right)$ is analytically intractable because the summation lies inside the logarithm.

\begin{theorem}[The Expectation-Maximization (EM) Algorithm for GMMs]
The EM algorithm alternates between two steps until convergence:
\begin{enumerate}
    \item \textbf{E-step (Expectation):} Compute the posterior probability (responsibility) that sample $\mathbf{x}_i$ was generated by component $k$ using current parameters $\boldsymbol{\theta}^{(t)}$:
    \begin{equation}
    \gamma_{ik} \triangleq P(z_i = k \mid \mathbf{x}_i, \boldsymbol{\theta}^{(t)}) = \frac{\pi_k^{(t)} \mathcal{N}(\mathbf{x}_i; \boldsymbol{\mu}_k^{(t)}, \Sigma_k^{(t)})}{\sum_{j=1}^K \pi_j^{(t)} \mathcal{N}(\mathbf{x}_i; \boldsymbol{\mu}_j^{(t)}, \Sigma_j^{(t)})}
    \label{eq:em_estep}
    \end{equation}
    \item \textbf{M-step (Maximization):} Update parameters to maximize the expected complete-data log-likelihood:
    \begin{align}
    N_k &= \sum_{i=1}^N \gamma_{ik} \label{eq:em_nk} \\
    \pi_k^{(t+1)} &= \frac{N_k}{N} \label{eq:em_pi} \\
    \boldsymbol{\mu}_k^{(t+1)} &= \frac{1}{N_k} \sum_{i=1}^N \gamma_{ik} \mathbf{x}_i \label{eq:em_mu} \\
    \Sigma_k^{(t+1)} &= \frac{1}{N_k} \sum_{i=1}^N \gamma_{ik} (\mathbf{x}_i - \boldsymbol{\mu}_k^{(t+1)})(\mathbf{x}_i - \boldsymbol{\mu}_k^{(t+1)})^T \label{eq:em_sigma}
    \end{align}
\end{enumerate}
\end{theorem}

\subsection{Fundamental Concentration Inequalities}
Concentration inequalities quantify the probability that a random variable deviates from its expected value.

\begin{theorem}[Markov's Inequality]
Let $X$ be a non-negative random variable ($X \ge 0$). For any $a > 0$:
\begin{equation}
P(X \ge a) \le \frac{\E[X]}{a}
\label{eq:markov_ineq}
\end{equation}
\end{theorem}

\begin{proof}
Let $\mathbb{I}(X \ge a)$ be the indicator variable for the event $X \ge a$. Because $X \ge 0$:
\begin{equation}
a \, \mathbb{I}(X \ge a) \le X
\end{equation}
Taking expectations on both sides:
\begin{equation}
a \, \E[\mathbb{I}(X \ge a)] = a P(X \ge a) \le \E[X] \implies P(X \ge a) \le \frac{\E[X]}{a}
\end{equation}
\end{proof}

\begin{theorem}[Chebyshev's Inequality]
Let $X$ have finite mean $\mu = \E[X]$ and variance $\sigma^2 = \Var(X)$. For any $\epsilon > 0$:
\begin{equation}
P(|X - \mu| \ge \epsilon) \le \frac{\sigma^2}{\epsilon^2}
\label{eq:chebyshev_ineq}
\end{equation}
\end{theorem}

\begin{proof}
Apply Markov's inequality to the non-negative random variable $Y = (X - \mu)^2$ with $a = \epsilon^2$:
\begin{equation}
P(|X - \mu| \ge \epsilon) = P((X - \mu)^2 \ge \epsilon^2) \le \frac{\E[(X - \mu)^2]}{\epsilon^2} = \frac{\sigma^2}{\epsilon^2}
\end{equation}
\end{proof}

\begin{theorem}[Hoeffding's Inequality]
Let $X_1, X_2, \dots, X_N$ be independent bounded random variables such that $X_i \in [a_i, b_i]$ with probability 1. Let $\bar{X} = \frac{1}{N}\sum_{i=1}^N X_i$ and $\mu = \E[\bar{X}]$. For any $\epsilon > 0$:
\begin{equation}
P(|\bar{X} - \mu| \ge \epsilon) \le 2 \exp\left( -\frac{2N^2 \epsilon^2}{\sum_{i=1}^N (b_i - a_i)^2} \right)
\label{eq:hoeffding_ineq}
\end{equation}
If all $X_i \in [0, 1]$, this simplifies to:
\begin{equation}
P(|\bar{X} - \mu| \ge \epsilon) \le 2 \exp\left( -2N\epsilon^2 \right)
\end{equation}
\end{theorem}

% ----------------------------------------------------------------------
% PART 3: VISUAL & GEOMETRIC INTUITION
% ----------------------------------------------------------------------
\section{Visual \& Geometric Intuition: The Geometry of Gaussian Mixture Density}

Figure~\ref{fig:gmm_geometry} depicts a two-component Gaussian Mixture Model in $\mathbb{R}^2$, showing the elliptical contour lines determined by the covariance matrices $\Sigma_1, \Sigma_2$ and the soft probabilistic responsibility boundary between clusters.

\begin{figure}[htbp]
\centering
\begin{tikzpicture}[
    scale=0.95,
    font=\sffamily\footnotesize
]

% Component 1: Center (-2.0, 0.5)
\begin{scope}[shift={(-2.0, 0.5)}, rotate=25]
    \draw[thick, color=tealborder!40] (0, 0) ellipse (3.0cm and 1.5cm);
    \draw[thick, color=tealborder!70] (0, 0) ellipse (2.0cm and 1.0cm);
    \draw[thick, color=tealborder] (0, 0) ellipse (1.0cm and 0.5cm);
    \filldraw[color=tealborder] (0, 0) circle (2.5pt) node[below] {$\boldsymbol{\mu}_1$};
\end{scope}
\node[above, color=tealborder] at (-2.0, 2.3) {\textbf{Cluster 1: } $\mathcal{N}(\boldsymbol{\mu}_1, \Sigma_1)$};

% Component 2: Center (2.5, -0.5)
\begin{scope}[shift={(2.5, -0.5)}, rotate=-35]
    \draw[thick, color=coralborder!40] (0, 0) ellipse (2.8cm and 1.3cm);
    \draw[thick, color=coralborder!70] (0, 0) ellipse (1.8cm and 0.8cm);
    \draw[thick, color=coralborder] (0, 0) ellipse (0.9cm and 0.4cm);
    \filldraw[color=coralborder] (0, 0) circle (2.5pt) node[below] {$\boldsymbol{\mu}_2$};
\end{scope}
\node[above, color=coralborder] at (2.5, 1.2) {\textbf{Cluster 2: } $\mathcal{N}(\boldsymbol{\mu}_2, \Sigma_2)$};

% Soft decision boundary (gamma_1 = gamma_2 = 0.5)
\draw[dashed, line width=1.5pt, color=primarydark] (0.2, 3.0) -- (0.5, -3.0);
\node[below, color=primarydark] at (0.5, -3.0) {Equi-Responsibility: $\gamma_{i1} = \gamma_{i2} = 0.5$};

% Intermediate sample point with soft assignment
\coordinate (X) at (0.0, 0.8);
\filldraw[color=purpleex] (X) circle (3pt) node[above right] {$\mathbf{x}_i$};
\node[draw=purpleborder, fill=purplebg, rounded corners=2pt, align=center] at (-0.2, -1.8) {
    \textbf{Soft Responsibility:}\\$\gamma_{i1} = 0.68, \; \gamma_{i2} = 0.32$
};

\end{tikzpicture}
\caption{The Geometry of Gaussian Mixture Models: Two distinct Gaussian clusters in $\mathbb{R}^2$ with overlapping contours, generating continuous soft posterior responsibilities $\gamma_{ik}$.}
\label{fig:gmm_geometry}
\end{figure}

% ----------------------------------------------------------------------
% PART 4: DEEP INTUITION & CONCEPTUAL DERIVATIONS
% ----------------------------------------------------------------------
\section{Deep Intuition \& Conceptual Derivations}

\subsection{Derivation of the EM Algorithm via the Evidence Lower Bound (ELBO)}
Let $\mathbf{x}$ be observed data and $\mathbf{z}$ be latent variables. The marginal log-likelihood is $\log p(\mathbf{x}; \boldsymbol{\theta})$.
Let $q(\mathbf{z})$ be an arbitrary probability distribution over the latent variables $\mathbf{z}$.
Using Jensen's inequality on the strictly concave logarithm function:
\begin{align}
\log p(\mathbf{x}; \boldsymbol{\theta}) &= \log \sum_{\mathbf{z}} p(\mathbf{x}, \mathbf{z}; \boldsymbol{\theta}) = \log \sum_{\mathbf{z}} q(\mathbf{z}) \frac{p(\mathbf{x}, \mathbf{z}; \boldsymbol{\theta})}{q(\mathbf{z})} \nonumber \\
&= \log \E_{q} \left[ \frac{p(\mathbf{x}, \mathbf{z}; \boldsymbol{\theta})}{q(\mathbf{z})} \right] \ge \E_{q} \left[ \log \frac{p(\mathbf{x}, \mathbf{z}; \boldsymbol{\theta})}{q(\mathbf{z})} \right] \triangleq \text{ELBO}(q, \boldsymbol{\theta})
\label{eq:elbo_jensen}
\end{align}
Expanding the Evidence Lower Bound:
\begin{align}
\text{ELBO}(q, \boldsymbol{\theta}) &= \sum_{\mathbf{z}} q(\mathbf{z}) \log p(\mathbf{x}, \mathbf{z}; \boldsymbol{\theta}) - \sum_{\mathbf{z}} q(\mathbf{z}) \log q(\mathbf{z}) \nonumber \\
&= \sum_{\mathbf{z}} q(\mathbf{z}) \log \left( p(\mathbf{z} \mid \mathbf{x}; \boldsymbol{\theta}) p(\mathbf{x}; \boldsymbol{\theta}) \right) + \mathcal{H}(q) \nonumber \\
&= \log p(\mathbf{x}; \boldsymbol{\theta}) - D_{\text{KL}}\left( q(\mathbf{z}) \,\big\|\, p(\mathbf{z} \mid \mathbf{x}; \boldsymbol{\theta}) \right)
\end{align}
This decomposition reveals the exact mechanics of the EM algorithm:
\begin{itemize}
    \item \textbf{E-step:} Maximize $\text{ELBO}(q, \boldsymbol{\theta}^{(t)})$ with respect to $q(\mathbf{z})$ while holding parameters $\boldsymbol{\theta}^{(t)}$ fixed. The ELBO touches the true marginal log-likelihood when the KL divergence vanishes, which occurs if and only if:
    \begin{equation}
    q^*(\mathbf{z}) = p(\mathbf{z} \mid \mathbf{x}; \boldsymbol{\theta}^{(t)})
    \end{equation}
    This is precisely the posterior responsibility $\gamma_{ik}$ computed in the E-step!
    \item \textbf{M-step:} Maximize $\text{ELBO}(q^*, \boldsymbol{\theta})$ with respect to parameters $\boldsymbol{\theta}$ while holding $q^*$ fixed.
    Because $\text{ELBO}(q^*, \boldsymbol{\theta}^{(t+1)}) \ge \text{ELBO}(q^*, \boldsymbol{\theta}^{(t)}) = \log p(\mathbf{x}; \boldsymbol{\theta}^{(t)})$, this guarantees that the true marginal log-likelihood increases monotonically at every iteration:
    \begin{equation}
    \log p(\mathbf{x}; \boldsymbol{\theta}^{(t+1)}) \ge \log p(\mathbf{x}; \boldsymbol{\theta}^{(t)})
    \end{equation}
\end{itemize}

\subsection{Connection Between GMM and $K$-Means Clustering}
$K$-Means clustering is mathematically equivalent to a special degenerate limit of a Gaussian Mixture Model:
\begin{enumerate}
    \item Assume all $K$ covariance matrices are isotropic and identical: $\Sigma_k = \sigma^2 I_d$.
    \item Assume uniform mixture priors: $\pi_k = \frac{1}{K}$.
    \item In the limit as variance $\sigma^2 \to 0$ (zero temperature limit):
    \begin{equation}
    \lim_{\sigma^2 \to 0} \gamma_{ik} = \lim_{\sigma^2 \to 0} \frac{\exp\left( -\frac{\|\mathbf{x}_i - \boldsymbol{\mu}_k\|^2}{2\sigma^2} \right)}{\sum_{j=1}^K \exp\left( -\frac{\|\mathbf{x}_i - \boldsymbol{\mu}_j\|^2}{2\sigma^2} \right)} = \begin{cases} 1 & \text{if } k = \argmin_j \|\mathbf{x}_i - \boldsymbol{\mu}_j\|^2 \\ 0 & \text{otherwise} \end{cases}
    \end{equation}
\end{enumerate}
Soft responsibilities collapse into hard binary cluster assignments, and the M-step mean update becomes the exact cluster centroid calculation of $K$-Means.

% ----------------------------------------------------------------------
% PART 5: EXPLICIT GOTCHAS & DIAGNOSTIC MISCONCEPTIONS
% ----------------------------------------------------------------------
\section{Explicit Gotchas \& Diagnostic Misconceptions}

\begin{trapbox}{Gotcha 1: Covariance Collapse / Likelihood Singularity in GMMs}
During EM training, a Gaussian component's mean $\boldsymbol{\mu}_k$ lands exactly on an isolated sample $\mathbf{x}_i$. The solver attempts to update $\Sigma_k$, causing $\det(\Sigma_k) \to 0$ and likelihood to explode to $+\infty$.
\textbf{The Trap:} The maximum likelihood formulation for GMMs has **unbounded singularities** on the boundary of parameter space. If a component collapses onto a single point, $\sigma^2 \to 0 \implies \frac{1}{\sigma^d} \to \infty$. Practical implementations must regularize the covariance matrices by adding a small positive ridge diagonal: $\Sigma_k \leftarrow \Sigma_k + \epsilon I_d$ (where $\epsilon \approx 10^{-6}$).
\end{trapbox}

\begin{trapbox}{Gotcha 2: Believing the EM Algorithm Finds the Global Maximum}
An analyst runs EM once on a complex dataset and assumes the parameters represent the global maximum of the marginal log-likelihood.
\textbf{The Trap:} The EM algorithm is a **local hill-climbing procedure**. The mixture log-likelihood surface is non-convex and filled with numerous suboptimal local maxima and saddle points. GMM implementations must be run multiple times with different random restarts (often initialized using $K$-Means centroids) to select the solution with the highest converged log-likelihood.
\end{trapbox}

\begin{trapbox}{Gotcha 3: Applying Markov's Inequality to Variables with Negative Support}
A student applies Markov's inequality to standard normal variable $Z \sim \mathcal{N}(0, 1)$ to bound $P(Z \ge 2) \le \frac{\E[Z]}{2} = 0$.
\textbf{The Trap:} Markov's inequality requires the random variable to be strictly **non-negative** ($X \ge 0$). Applying it to variables that take negative values produces nonsensical probabilities (such as upper-bounding a non-zero probability by zero or a negative number).
\end{trapbox}

\begin{trapbox}{Gotcha 4: Conflating Sample Variance with Variance of the Sample Mean}
A practitioner samples $N = 100$ points from a distribution with variance $\sigma^2 = 16$. They assert that the sample mean $\bar{X}$ has standard deviation $4$.
\textbf{The Trap:} The individual samples have variance $\sigma^2 = 16$ (standard deviation $4$). However, by independence, the **sample mean** has variance $\Var(\bar{X}) = \frac{\sigma^2}{N} = \frac{16}{100} = 0.16$, yielding a standard error of $\frac{\sigma}{\sqrt{N}} = \frac{4}{10} = 0.4$. The sample mean concentrates around $\mu$ at rate $\mathcal{O}(1/\sqrt{N})$.
\end{trapbox}

% ----------------------------------------------------------------------
% PART 6: COMPARATIVE SYNTHESIS & SUMMARY TABLE
% ----------------------------------------------------------------------
\section{Comparative Synthesis \& Summary Table}

Table~\ref{tab:clustering_comp} compares the characteristics of clustering and latent density models.

\begin{table}[htbp]
\centering
\small
\begin{tabularx}{\textwidth}{lXXXXX}
\toprule
\textbf{Model} & \textbf{Cluster Assignment} & \textbf{Cluster Geometry} & \textbf{Objective Function} & \textbf{Computational Complexity} & \textbf{Density Output} \\
\midrule
$K$-Means & Hard (Binary $0/1$) & Spherical (isotropic Voronoi cells) & Residual sum of squares $\sum \|\mathbf{x}_i - \boldsymbol{\mu}_{z_i}\|^2$ & $\mathcal{O}(N K d \cdot \text{iters})$ & No (Centroids only) \\
\addlinespace
GMM (EM) & Soft (Probabilities $\gamma_{ik} \in [0, 1]$) & Ellipsoidal (arbitrary full covariance $\Sigma_k$) & Marginal log-likelihood $\sum \log \sum \pi_k \mathcal{N}(\mathbf{x}_i)$ & $\mathcal{O}(N K d^2 \cdot \text{iters})$ & Yes (Complete generative density $p(\mathbf{x})$) \\
\addlinespace
Factor Analysis & Continuous latent $\mathbf{z} \in \mathbb{R}^k$ & Linear manifold + diagonal noise & Marginal Gaussian log-likelihood & $\mathcal{O}(N d k \cdot \text{iters})$ & Yes (Low-rank Gaussian density) \\
\addlinespace
KDE & Non-parametric & Kernel balls centered on every data point & Integrated squared error / leave-one-out likelihood & $\mathcal{O}(N^2 d)$ at inference & Yes (Continuous empirical density) \\
\bottomrule
\end{tabularx}
\caption{Comparative Synthesis: Clustering and Density Estimation Paradigms.}
\label{tab:clustering_comp}
\end{table}

% ----------------------------------------------------------------------
% PART 7: PYTHON / SCIENTIFIC COMPUTING LAB
% ----------------------------------------------------------------------
\section{Python / Scientific Computing Lab}

The following production script implements a complete Gaussian Mixture Model with the Expectation-Maximization algorithm from scratch, verifying the monotonic ascent of the marginal log-likelihood.

\begin{lstlisting}[language=Python, caption={Gaussian Mixture Model EM Algorithm Implementation in Python.}]
import numpy as np

class GaussianMixtureEM:
    def __init__(self, n_components: int, max_iter: int = 100, tol: float = 1e-4, seed: int = 42):
        self.K = n_components
        self.max_iter = max_iter
        self.tol = tol
        self.seed = seed
        self.weights_ = None
        self.means_ = None
        self.covariances_ = None
        self.log_likelihood_history_ = []

    def _multivariate_gaussian_log_pdf(self, X: np.ndarray, mean: np.ndarray, cov: np.ndarray) -> np.ndarray:
        N, d = X.shape
        # Regularize covariance to prevent collapse
        cov_reg = cov + 1e-6 * np.eye(d)
        diff = X - mean
        sign, logdet = np.linalg.slogdet(cov_reg)
        inv_cov = np.linalg.inv(cov_reg)
        mahalanobis = np.sum(diff @ inv_cov * diff, axis=1)
        log_prob = -0.5 * (d * np.log(2.0 * np.pi) + logdet + mahalanobis)
        return log_prob

    def fit(self, X: np.ndarray):
        N, d = X.shape
        rng = np.random.RandomState(self.seed)

        # Initialize parameters: random points as initial means
        self.weights_ = np.ones(self.K) / self.K
        random_idx = rng.choice(N, self.K, replace=False)
        self.means_ = X[random_idx].copy()
        self.covariances_ = np.array([np.cov(X, rowvar=False) for _ in range(self.K)])

        for it in range(self.max_iter):
            # E-step: compute log responsibilities via log-sum-exp
            log_weighted_probs = np.zeros((N, self.K))
            for k in range(self.K):
                log_weighted_probs[:, k] = np.log(self.weights_[k] + 1e-12) + \
                                          self._multivariate_gaussian_log_pdf(X, self.means_[k], self.covariances_[k])

            # Total log-likelihood
            max_log = np.max(log_weighted_probs, axis=1, keepdims=True)
            log_sum_exp = max_log + np.log(np.sum(np.exp(log_weighted_probs - max_log), axis=1, keepdims=True))
            current_log_lik = np.sum(log_sum_exp)
            self.log_likelihood_history_.append(current_log_lik)

            # Normalized responsibilities: gamma_ik = exp(log_weighted - log_sum_exp)
            gamma = np.exp(log_weighted_probs - log_sum_exp)

            # M-step: parameter updates
            N_k = np.sum(gamma, axis=0) + 1e-12
            self.weights_ = N_k / N

            for k in range(self.K):
                self.means_[k] = np.sum(gamma[:, k:k+1] * X, axis=0) / N_k[k]
                diff = X - self.means_[k]
                self.covariances_[k] = (diff.T @ (gamma[:, k:k+1] * diff)) / N_k[k]

            # Check convergence
            if it > 0 and abs(self.log_likelihood_history_[-1] - self.log_likelihood_history_[-2]) < self.tol:
                break

        return self

def main():
    rng = np.random.RandomState(101)
    N = 300
    # Generate 2-component synthetic data
    X1 = rng.randn(N // 2, 2) + np.array([3.0, 3.0])
    X2 = rng.randn(N // 2, 2) + np.array([-2.0, -1.0])
    X = np.vstack([X1, X2])

    gmm = GaussianMixtureEM(n_components=2, max_iter=50)
    gmm.fit(X)

    print("Gaussian Mixture Model EM Execution Complete:")
    print(f"Number of Iterations: {len(gmm.log_likelihood_history_)}")
    print(f"Initial Log-Likelihood: {gmm.log_likelihood_history_[0]:.2f}")
    print(f"Final Log-Likelihood:   {gmm.log_likelihood_history_[-1]:.2f}")

    # Verify monotonic ascent
    diffs = np.diff(gmm.log_likelihood_history_)
    is_monotonic = np.all(diffs >= -1e-6)
    print(f"Log-Likelihood Monotonically Ascending: {is_monotonic}")
    print("Estimated Component Means:")
    for k, mean in enumerate(gmm.means_):
        print(f"  Cluster {k+1}: [{mean[0]:.3f}, {mean[1]:.3f}], Weight: {gmm.weights_[k]:.3f}")

if __name__ == "__main__":
    main()
\end{lstlisting}

% ----------------------------------------------------------------------
% PART 8: COMPREHENSIVE PRACTICE PROBLEMS
% ----------------------------------------------------------------------
\section{Comprehensive Practice Problems}

\subsection*{Part A: Analytical \& Mathematical Problems}

\begin{enumerate}[label=\textbf{P\arabic*.}]
    \item \textbf{Analytical 2D Mahalanobis Distance and Ellipsoidal Contours:}
    Let $\mathbf{X} \sim \mathcal{N}(\boldsymbol{\mu}, \Sigma)$ with $\boldsymbol{\mu} = \begin{bmatrix} 1 \\ 2 \end{bmatrix}$ and $\Sigma = \begin{bmatrix} 4 & 2 \\ 2 & 2 \end{bmatrix}$.
    \begin{enumerate}[label=(\alph*)]
        \item Compute the determinant $|\Sigma|$ and inverse matrix $\Sigma^{-1}$.
        \item Compute the Mahalanobis distance $D_M(\mathbf{x}, \boldsymbol{\mu})$ for observation $\mathbf{x} = [3, 4]^T$.
        \item Find the eigenvalues and eigenvectors of $\Sigma$ and describe the principal orientation and axis lengths of the $1$-standard-deviation contour ellipsoid.
    \end{enumerate}

    \item \textbf{Unbiasedness and Degrees of Freedom of Sample Covariance:}
    Let $\mathbf{X}_1, \dots, \mathbf{X}_N \stackrel{\text{i.i.d.}}{\sim} \mathcal{N}(\boldsymbol{\mu}, \Sigma)$.
    Prove that the Maximum Likelihood Estimator $\hat{\Sigma}_{\text{MLE}} = \frac{1}{N}\sum_{i=1}^N (\mathbf{X}_i - \bar{\mathbf{X}})(\mathbf{X}_i - \bar{\mathbf{X}})^T$ has expectation:
    \begin{equation}
    \E[\hat{\Sigma}_{\text{MLE}}] = \frac{N - 1}{N} \Sigma
    \end{equation}

    \item \textbf{Analytical E-step Calculation for a 1D GMM:}
    Consider a 1D GMM with $K = 2$ components: $\pi_1 = 0.6, \mu_1 = 0, \sigma_1^2 = 1$ and $\pi_2 = 0.4, \mu_2 = 3, \sigma_2^2 = 1$.
    For an observation $x = 1.5$:
    \begin{enumerate}[label=(\alph*)]
        \item Evaluate the Gaussian densities $\mathcal{N}(1.5; 0, 1)$ and $\mathcal{N}(1.5; 3, 1)$.
        \item Compute the posterior responsibilities $\gamma_1$ and $\gamma_2$.
        \item State which cluster $x = 1.5$ is more likely to belong to, and explain why despite being equidistant from both means.
    \end{enumerate}

    \item \textbf{Sample Complexity via Hoeffding's Inequality:}
    A pollster wants to estimate the true proportion $p \in [0, 1]$ of voters favoring a policy using empirical average $\bar{X}_N$.
    Using Hoeffding's inequality, calculate the minimum number of independent samples $N$ required to guarantee that $|\bar{X}_N - p| \le 0.02$ with confidence at least $95\%$ ($\delta = 0.05$).

    \item \textbf{Chebyshev vs. Markov Bound Comparison:}
    Let $X$ be the number of heads in $N = 100$ independent flips of a fair coin.
    \begin{enumerate}[label=(\alph*)]
        \item Find the exact mean $\mu = \E[X]$ and variance $\sigma^2 = \Var(X)$.
        \item Use Markov's inequality to bound $P(X \ge 75)$.
        \item Use Chebyshev's inequality to bound $P(X \ge 75) \le P(|X - 50| \ge 25)$.
        \item Compare the bounds and explain why Chebyshev is dramatically tighter.
    \end{enumerate}
\end{enumerate}

\subsection*{Part B: Diagnostic Multiple-Choice Questions}

\begin{enumerate}[label=\textbf{MCQ\arabic*.}]
    \item What is guaranteed regarding the marginal log-likelihood $\sum_{i=1}^N \log p(\mathbf{x}_i; \boldsymbol{\theta})$ at each successive iteration of the Expectation-Maximization (EM) algorithm?
    \begin{enumerate}[label=(\alph*)]
        \item It is guaranteed to decrease monotonically toward zero.
        \item It is guaranteed to increase monotonically (or remain constant) at every iteration: $\ell(\boldsymbol{\theta}^{(t+1)}) \ge \ell(\boldsymbol{\theta}^{(t)})$.
        \item It is guaranteed to converge to the global maximum of the parameter space.
        \item It evaluates to a strictly negative value if and only if the data is linearly separable.
    \end{enumerate}

    \item If a bounded random variable $X_i \in [0, 1]$ has unknown distribution with mean $\mu$, which concentration inequality provides an exponential bound in sample size $N$ on the deviation of sample mean $|\bar{X}_N - \mu|$?
    \begin{enumerate}[label=(\alph*)]
        \item Markov's Inequality
        \item Chebyshev's Inequality
        \item Hoeffding's Inequality
        \item Cauchy-Schwarz Inequality
    \end{enumerate}
\end{enumerate}

